\section{Роль расчёта термодинамических свойств воды в инженерных задачах}

Вода и водяной пар являются базовой рабочей средой для теплоэнергетики,
химико-технологических процессов и смежных инженерных систем.
На этапах проектирования, наладки и эксплуатации требуется многократный пересчёт
состояния по ограниченному набору измеряемых параметров.
К типовым задачам относятся расчёты тепловых схем, выбор режимов котлоагрегатов и турбин,
оценка потерь в трубопроводах, анализ теплообмена и переходных процессов.

Результат расчёта используется не только как набор чисел.
Он служит основой для диагностических диаграмм и принятия эксплуатационных решений.
Диаграммы $T$--$s$, $p$--$h$, $p$--$v$ позволяют контролировать положение точки
относительно линии насыщения и быстрее выявлять недопустимые режимы.

С инженерной точки зрения критичны три свойства вычислительного инструмента.
Первое свойство --- скорость, достаточная для интерактивной работы и пакетной обработки.
Второе свойство --- устойчивость на фазовых границах и в околокритической области.
Третье свойство --- диагностируемость: при ошибочном вводе система должна выдавать
явную причину отказа, а не неинтерпретируемый результат.

\section{Стандарты и существующие подходы к расчёту свойств воды и пара}

Исторически расчёты выполнялись по паровым таблицам и диаграммам с интерполяцией~\cite{alexandrov_grigorev_1999}.
Современные программные реализации опираются на аналитические зависимости
и дополняют их численными методами для обратных маршрутов.

В инженерной практике применяются стандартизованные формулировки IAPWS.
Они обеспечивают сопоставимость результатов, но требуют строгого соблюдения
границ применимости и правил перехода между областями состояния.
Ключевая особенность предметной области --- региональная структура модели.
Разные участки рабочей области описываются разными уравнениями состояния
и разными независимыми переменными.

Наиболее сложным с вычислительной точки зрения является регион~3.
Для этой области естественными переменными являются плотность и температура,
поэтому при других входных парах требуются обратные зависимости и численные решатели.
Следовательно, практическая программная реализация должна включать не только формулы,
но и маршрутизацию расчёта, валидацию входа и обработку граничных случаев.

\subsection{Эволюция стандартов: от IFC-67 к IAPWS-IF97}

Для промышленного применения свойства воды и водяного пара
долгое время рассчитывались по формулировке IFC-67,
которая была положена в основу изданий ASME Steam Tables 1967 года
и многочисленных национальных таблиц и справочников~\cite{iapws_faq_asme}.
Именно эта традиция табличного представления расчётных данных
сохранилась и в российской инженерной практике,
что подтверждается поздними русскоязычными справочными изданиями
и учебной литературой~\cite{alexandrov_grigorev_1999,gidaspov_severina_2014}.

Однако развитие вычислительной техники
и рост требований к численной скорости расчёта
сделали табличный подход недостаточным.
В конце 1997 года IAPWS приняла промышленную формулировку IAPWS-IF97,
которая официально заменила IFC-67
и стала международным стандартом для расчётов в паросиловой отрасли~\cite{iapws_faq_asme}.
Новая формулировка сохраняет инженерную направленность,
но строится уже как набор региональных аналитических уравнений,
ориентированных на быструю и устойчивую машинную обработку.

По сравнению с IAPWS-95,
которая предназначена для общего и научного применения,
IF97 сознательно оптимизирована под промышленную эксплуатацию:
она быстрее вычисляется,
дольше сохраняется без смены нормативной базы
и лучше соответствует задачам расчёта тепловых схем,
оборудования и режимов работы энергетических установок~\cite{iapws_faq_asme,coolprop_if97_docs}.
Для дипломной работы это важно по двум причинам.
Во-первых, сама постановка задачи относится к инженерным, а не к фундаментальным расчётам.
Во-вторых, требуется единое ядро,
пригодное как для пользовательских интерфейсов,
так и для сервисной пакетной обработки.

\subsection{Сравнительный анализ существующих библиотек}

Современная программная практика уже предлагает готовые библиотеки,
использующие IF97 или совместимые с ней маршруты расчёта.
Наиболее показательными для сравнения являются библиотека CoolProp
и специализированная библиотека SteamTables.jl.

CoolProp представляет собой зрелую многокомпонентную библиотеку,
ориентированную на широкий круг термодинамических расчётов.
Для воды и пара в ней предусмотрена отдельная реализация \texttt{IF97::Water},
основанная на последних релизах IAPWS по термодинамическим и транспортным свойствам~\cite{coolprop_if97_docs}.
Преимуществами CoolProp являются
широкий набор обёрток для разных языков,
большое число поддерживаемых выходных свойств
и наличие большинства типовых входных пар.
При этом IF97-режим в CoolProp остаётся водо-паровым частным случаем:
он ограничен более узкой областью применимости,
не поддерживает смеси
и не реализует универсальные частные производные в общем интерфейсе~\cite{coolprop_if97_docs}.

SteamTables.jl, напротив, является компактной предметной библиотекой,
сосредоточенной только на воде и водяном паре.
Она реализует свойства IF97,
поддерживает входные постановки $(p, T)$, $(p, h)$ и $(p, s)$,
а также предоставляет функции насыщения
и возможность использовать физические единицы через \texttt{Unitful.jl}~\cite{steamtables_jl}.
Такой подход удобен для расчётных сценариев в среде Julia,
но библиотека ориентирована на более узкий круг маршрутов
и на один язык программирования.

В Python-среде наиболее распространённым решением остаётся пакет \texttt{iapws}.
Он удобен для учебных расчётов, верификации формул и быстрых скриптов,
но его сильная сторона проявляется именно в простоте использования,
а не в производительности массовой обработки.
При переходе к большому числу запросов начинают сказываться ограничения
интерпретируемого окружения и накладные расходы на вызовы функций.

В прикладных системах также нередко встречаются REST-обёртки над термодинамическими
библиотеками, в том числе над CoolProp.
Такой вариант упрощает интеграцию с внешними клиентами и веб-сервисами,
но добавляет сетевые задержки, сериализацию и дополнительные преобразования данных.
Поэтому сравнение на рисунке~\ref{fig:expected_performance}
следует воспринимать как сопоставление типовых классов решений,
а не как микробенчмарк в идентичных условиях исполнения.

Ниже приведена гистограмма, которая суммирует это сравнение:
Python-пакет \texttt{iapws}, REST-представление CoolProp,
REST-реализацию на Rust и ожидаемые показатели разрабатываемого комплекса.
Для собственного комплекса здесь показана консервативная проектная оценка,
а не итоговый измеренный результат.
Логарифмическая шкала выбрана намеренно,
поскольку разброс между сценариями составляет несколько порядков.

\begin{figure}[htbp]
    \centering
    \resizebox{\linewidth}{!}{
    \begin{tikzpicture}
    \begin{axis}[
        ybar,
        ymode=log,
        ymajorgrids=true,
        grid style={dashed, gray!30},
        ylabel={Производительность, точек/с},
        symbolic x coords={PY,CP,RR,RG},
        xtick=data,
        xticklabels={
            {Python\\(iapws)},
            {CoolProp\\(REST)},
            {Rust\\(REST)},
            {Наст. работа\\(ожидаемое)}
        },
        xticklabel style={align=center, font=\footnotesize},
        nodes near coords,
        nodes near coords style={
            font=\tiny,
            /pgf/number format/fixed,
            /pgf/number format/1000 sep={\,},
            anchor=south
        },
        bar width=0.8cm,
        enlarge x limits=0.2,
        ymin=500,
        width=\linewidth,
        height=7cm
    ]
        % Аналоги (сплошная заливка)
        \addplot[fill=gray!30, draw=black] coordinates {
            (PY, 1500)
            (CP, 80000)
            (RR, 320000)
        };
        % Твоя система (полупрозрачная)
        \addplot[fill=blue!40, draw=blue, opacity=0.4] coordinates {
            (RG, 1300000)
        };
    \end{axis}
    \end{tikzpicture}
    }
    \caption{Сравнительный анализ производительности аналогов и ожидаемые показатели разрабатываемого комплекса (проектные значения)}
    \label{fig:expected_performance}
\end{figure}

Сопоставление этих решений показывает,
что существующие библиотеки хорошо закрывают задачу одиночного вызова формул,
но не решают автоматически архитектурную задачу проекта:
единое вычислительное ядро на Rust,
общее DTO-представление для разных клиентов,
маршрут $\rho T$ на всей рабочей области,
контролируемую маршрутизацию по регионам
и повторное использование одного и того же ядра
в настольном интерфейсе, оболочке на базе веб-технологий и gRPC-сервисе.
Следовательно, в рамках дипломной работы требуется
не только выбрать стандарт расчёта,
но и реализовать программный комплекс,
в котором физическая модель, архитектура и эксплуатационные требования
согласованы между собой.
